%% ma: L12-B14-0120
%% lop: 12
%% chuong: 5
%% bai: 14
%% dang: 12.14.34
%% loai: DS
%% muc: [NB, TH, TH, VD]
%% dap_an: "ĐSĐS"
%% nguon: "Ảnh giáo viên gửi 07/11/2026 (Tư Duy Mở - tuduymo.com, đề tự luyện Phương trình mặt phẳng, đề khoảng cách), Câu 16"
%% kien_thuc: "Bài 14 (hai mặt phẳng song song, khoảng cách từ một điểm đến mặt phẳng)"
%% trang_thai: chinh-thuc
%% ghi_chu: "Tài liệu gốc không kèm đáp án; đáp án do Claude giải và kiểm bằng sympy. Dạng đã duyệt 07/11/2026."
\begin{ex}
Trong không gian với hệ trục tọa độ $Oxyz$, cho mặt phẳng $(P)\colon x+3y+z-3=0$ và mặt phẳng $(Q)\colon -x-3y-z-11=0$.
\choiceTF
{\True $(P)\parallel(Q)$.}
{$d\left((P),(Q)\right)=8$.}
{\True Tập hợp các điểm cách đều $(P)$, $(Q)$ là mặt phẳng $(R)\colon x+3y+z+4=0$.}
{Tập hợp các điểm nằm giữa $(P)$, $(Q)$ sao cho khoảng cách từ nó đến $(P)$ gấp $3$ lần khoảng cách từ nó đến $(Q)$ là mặt phẳng $(\alpha)\colon x+3y+z+3=0$.}
\loigiai{$(Q)\colon x+3y+z+11=0$. Đặt $f=x+3y+z$.

a) Cùng vectơ pháp tuyến, $-3\ne11$. Đúng.

b) $d=\dfrac{|-3-11|}{\sqrt{11}}=\dfrac{14}{\sqrt{11}}\ne8$. Sai.

c) $|f-3|=|f+11|\Leftrightarrow f=-4$. Đúng.

d) Điểm nằm giữa: $-11<f<3$, $3-f=3(f+11)\Leftrightarrow f=-\dfrac{15}2$: $(\alpha)\colon 2x+6y+2z+15=0$. Sai.}
\end{ex}
